% !PRINTED 186
\weberpage{186}

Durch Integration von (2) findet man:
\[
\sigma'(u) = \left(\frac{1}{u} - \frac{a_1 u^3}{3} - \frac{a_2 u^5}{5} - \frac{a_3 u^7}{7}\cdots\right)\sigma(u),
\]
und aus dieser Gleichung folgt:

% [sic] The source prints g_3^2 in both \alpha_3 below and the u^9 term of (7); homogeneity and (5) require g_2^2.
\weq{6}{\alpha = 0, \quad \alpha_1 = \frac{-g_2}{240}, \quad \alpha_2 = \frac{-g_3}{840}, \quad \alpha_3 = \frac{-g_3^2}{161\,280}.}

% [sic] The source prints \sigma'''(u)=0; the coefficient argument requires \sigma'''(0)=0, since \sigma''' is not identically zero.
Daß $\alpha=0$ ist, folgt auch aus dem in \S\,35 bewiesenen
$\sigma'''(u) = 0$; indessen war dies dort auf ziemlich umständlichem
Wege bewiesen. Wir haben sonach folgenden Anfang der Entwickelung
 von $\sigma(u)$, wobei der Übersicht halber die numerischen
Nenner in ihre Primfaktoren zerlegt sind:

\weq{7}{\sigma(u) = u - \frac{g_2}{2^4\cdot 3\cdot 5}u^5 - \frac{g_3}{2^3\cdot 3\cdot 5\cdot 7}u^7 - \frac{g_3^2}{2^9\cdot 3^2\cdot 5\cdot 7}u^9 - \cdots}

Zur rekurrenten Berechnung der Koeffizienten in der Potenzreihe
 für $\sigma(u)$ dient eine partielle Differentialgleichung, die als
eine Umformung der partiellen Differentialgleichung für die
$\vartheta$-Funktionen betrachtet werden kann. Wir leiten sie auf folgendem
 Wege her.

Die Funktion $\sigma(u)$ war durch die beiden Periodengleichungen
\S\,35, (9)

\weq{8}{
\begin{aligned}
&\sigma(u+\omega_1) = -e^{\eta_1(2u+\omega_1)}\sigma(u),\\
&\sigma(u+\omega_2) = -e^{\eta_2(2u+\omega_2)}\sigma(u),\\
&\eta_1\omega_2 - \eta_2\omega_1 = \pi i
\end{aligned}
}

bis auf einen von $u$ unabhängigen Faktor bestimmt. Es ergibt
sich aber durch Differentiation der letzten Gleichung (8):
\[
\begin{aligned}
\frac{\partial \eta_1}{\partial \omega_1}\omega_2 - \frac{\partial \eta_2}{\partial \omega_1}\omega_1 - \eta_2 &= 0,\\
\frac{\partial \eta_1}{\partial \omega_2}\omega_2 - \frac{\partial \eta_2}{\partial \omega_2}\omega_1 + \eta_1 &= 0,
\end{aligned}
\]
also:
\[
\omega_2\left(\eta_1\frac{\partial \eta_1}{\partial \omega_1} + \eta_2\frac{\partial \eta_1}{\partial \omega_2}\right) = \omega_1\left(\eta_1\frac{\partial \eta_2}{\partial \omega_1} + \eta_2\frac{\partial \eta_2}{\partial \omega_2}\right),
\]
und wir führen also eine Größe $r$ ein, die wir so definieren:

\weq{9}{
\begin{aligned}
&\eta_1\frac{\partial \eta_1}{\partial \omega_1} + \eta_2\frac{\partial \eta_1}{\partial \omega_1} = -r\omega_1, \\ % [sic] Second denominator is printed as partial omega_1; the derivation requires partial omega_2.
&\eta_1\frac{\partial \eta_1}{\partial \omega_2} + \eta_2\frac{\partial \eta_2}{\partial \omega_2} = -r\omega_2. % [sic] First fraction is printed partial eta_1 / partial omega_2; the derivation requires partial eta_2 / partial omega_1.
\end{aligned}
}
